\clearpage
\chapter*{Zusammenfassung}
Diese Arbeit untersucht, ob eine zusätzliche selbstüberwachte Anpassung bereits
vortrainierter Transformer den Bedarf an gelabelten Trainingsdaten für die
Phishing-Webseitenerkennung verringert und eine ressourcenschonende Verarbeitung
unterstützt. Auf PhreshPhish werden vier Repräsentationsvarianten verglichen:
unveränderte Basisencoder sowie Task-Adaptive Pretraining (TAPT) für URL,
HTML-Text oder beide Eingabemodalitäten. Lineare und nichtlineare Probes verwenden
sieben verschachtelte Labelbudgets von 2.000 bis 200.000 Webseiten.
Die finale Evaluation umfasst 168.060 Webseiten sowie vier überlappende Teilmengen
mit Domain-, Template- oder Zeitbedingungen. Schwellen werden auf einem getrennten
Bestand von 50.000 legitimen Webseiten für eine Ziel-False-Positive-Rate von 0,5\,\%
kalibriert.

Die kombinierte Anpassung RUT verbessert gegenüber unveränderten Encodern bei
20.000 Trainingslabels die mittlere True-Positive-Rate auf dem vollständigen Testbestand
um 6,08 Prozentpunkte (zehn gepaarte Läufe; 95-\%-Konfidenzintervall [5,38; 6,79]).
50.000 Labels bilden die kleinste untersuchte Budgetstufe, die die festgelegte
Nichtunterlegenheitsentscheidung gegenüber den mit 200.000 Labels trainierten
R0- und TF-IDF/XGBoost-Referenzen über alle fünf Bedingungen erfüllt. Die damit
verbundene Reduktion um 75\,\% betrifft die Downstream-Trainingslabels; Labels für
Stichprobenauswahl, Entwicklung und Kalibration sind darin nicht enthalten.

Auf dem externen Putra-Datensatz steigt die über 200 bis 2.000 Labels integrierte
Average-Precision-Lernkurve durch RUT um 2,54 Prozentpunkte. Eine weitere
MLM-Anpassung auf dem frühen Putra-Bestand verschlechtert diese Größe hingegen
um 1,14 Prozentpunkte. Der Nutzen zusätzlicher Anpassung ist somit korpus- und
protokollabhängig. Ein multimodaler Prototyp aus URL, HTML-Text und DOM erreicht
96,12\,\% TPR bei 0,666\,\% FPR. Seine URL-first-Kaskade erreicht 95,77\,\% TPR,
eskaliert 14,93\,\% der Testseiten und steigert den gemessenen In-Memory-Durchsatz
von 183,2 auf 447,0 Seiten pro Sekunde.

Die Ergebnisse stützen einen begrenzten Nutzen von TAPT für Label-Effizienz und
gestufte Verarbeitung. Feste TAPT-Checkpoints, nachträglich ergänzte Analysen,
überlappende Testbedingungen und die retrospektive externe Evaluation begrenzen
die Verallgemeinerbarkeit. Laufzeitmessungen erfassen den aufbereiteten Modellpfad
auf einer GPU und keine vollständige betriebliche Verarbeitungskette.

\clearpage
\begin{otherlanguage}{english}
\chapter*{Abstract}
This thesis examines whether additional self-supervised adaptation of pretrained
transformers reduces the labelled training data required for phishing website
detection and supports efficient processing. Four representation variants are
compared on PhreshPhish: unchanged base encoders and task-adaptive pretraining
(TAPT) of the URL encoder, the HTML text encoder, or both. Linear and nonlinear
probes use seven nested label budgets ranging from 2,000 to 200,000 websites.
Final evaluation covers 168,060 websites and four overlapping subsets defined by
domain, template, or time conditions. Thresholds are calibrated on a separate set
of 50,000 legitimate websites for a target false positive rate of 0.5\%.

At 20,000 training labels, the combined adaptation RUT improves mean true positive
rate over unchanged encoders by 6.08 percentage points on the complete test set
(ten paired runs; 95\% confidence interval [5.38, 6.79]). A budget of 50,000 labels
is the smallest tested level meeting the specified noninferiority decision against
both the R0 and TF-IDF/XGBoost references trained with 200,000 labels across all
five evaluation conditions. The resulting 75\% reduction concerns downstream
training labels; labels used for sample selection, development, and calibration
are excluded from this accounting.

On the external Putra dataset, RUT improves the integrated average precision
learning curve over budgets of 200 to 2,000 labels by 2.54 percentage points.
Additional masked language modelling on the early Putra data instead decreases
this measure by 1.14 percentage points. Benefits of further adaptation therefore
depend on the corpus and training protocol. A multimodal prototype combining
URL, HTML text, and DOM reaches a true positive rate of 96.12\% at a false positive
rate of 0.666\%. Its URL-first cascade reaches a true positive rate of 95.77\%,
escalates 14.93\% of test websites, and increases measured in-memory throughput
from 183.2 to 447.0 websites per second.

The results support a qualified benefit of TAPT for label efficiency and staged
processing. Fixed TAPT checkpoints, analyses added after initial experiments,
overlapping evaluation conditions, and retrospective external evaluation limit
generalisability. Runtime measurements cover the prepared model path on one GPU
and do not represent a complete operational processing pipeline.
\end{otherlanguage}
\clearpage
